\pagebreak
\section*{以往工作}
\label{sec:related}

开发通用算法长期以来一直是强化学习研究的目标。
PPO\citep{schulman2017ppo} 是使用最广泛的算法之一，相对稳健，但需要大量经验，而且性能往往低于专门方法。
SAC\citep{haarnoja2018sac} 是连续控制中的常用选择，并利用经验回放提升数据效率；但实践中仍需要调参，尤其是熵系数，并且在高维输入下表现困难\citep{yarats2019sacae}。
MuZero\citep{schrittwieser2019muzero} 使用价值预测模型进行规划，并已应用于棋盘游戏和 Atari；但作者没有发布实现，且算法包含多个复杂组件，因此复现有挑战。
Gato\citep{reed2022gato} 将一个大模型拟合到多任务专家示范上，但只有在存在专家数据时才适用。
相比之下，Dreamer 能在固定超参数下掌握多样化环境，不需要专家数据，并且实现已开源。

Minecraft 也是近期研究重点。Microsoft 通过 MALMO\citep{johnson2016malmo} 发布了该成功游戏的免费研究版本。
MineRL\citep{guss2019minerl} 提供多个竞赛环境，本文实验也以这些环境为基础。
MineRL 竞赛通过多样化人类数据集，支持智能体探索并学习有意义的技能\citep{guss2019minerl}。
Voyager\citep{wang2023voyager} 通过调用语言模型 API，在技术树中获得了与 Dreamer 相近深度的物品；但它运行在 MineFlayer 脚本层之上，该层是专门为游戏工程化设计的，并暴露高层动作。
VPT\citep{baker2022vpt} 基于承包者收集的键盘鼠标专家数据，通过行为克隆训练 Minecraft 智能体，并使用强化学习微调；其获得钻石需要 720 块 GPU 训练 9 天。
相比之下，Dreamer 使用 MineRL 竞赛动作空间，在没有人类数据的情况下，仅用 1 块 GPU 训练 9 天，就能从稀疏奖励中自主学会收集钻石。
